\documentclass[11pt,a4paper]{article}

\usepackage[T1]{fontenc}
\usepackage{lmodern}
\usepackage{amsmath,amssymb,amsthm,mathtools}
\usepackage{bm}
\usepackage{booktabs}
\usepackage[margin=29mm]{geometry}
\usepackage{microtype}
\usepackage[hidelinks]{hyperref}

\newtheorem{theorem}{Theorem}[section]
\newtheorem{proposition}[theorem]{Proposition}
\newtheorem{lemma}[theorem]{Lemma}
\newtheorem{corollary}[theorem]{Corollary}
\theoremstyle{definition}
\newtheorem{definition}[theorem]{Definition}
\theoremstyle{remark}
\newtheorem{remark}[theorem]{Remark}

\newcommand{\End}{\operatorname{End}}
\newcommand{\tr}{\operatorname{tr}}
\newcommand{\dd}{\mathrm{d}}
\newcommand{\supp}{\operatorname{supp}}
\newcommand{\cQ}{\mathcal{Q}}
\newcommand{\tcQ}{\widetilde{\mathcal{Q}}}
\newcommand{\cyc}{\mathsf{S}}

\title{Algebraic endpoint densities are the transfer-matrix charges\\
in regular Yang--Baxter chains}
\author{Xiaojian Zeng\\
\small PBC School of Finance, Tsinghua University, Beijing 100190, China\\
\small \texttt{xiaojianzeng@foxmail.com}}
\date{14 October 2026}

\begin{document}
\maketitle

\begin{abstract}
Schouten and de Leeuw recently introduced algebraic charge densities obtained by
differentiating directed products of a regular Yang--Baxter $R$-matrix and conjectured
that two endpoint differences represent the usual logarithmic-transfer charges.  We
prove this conjecture for every finite-dimensional local space, every regular
braiding-normalized (not necessarily difference-form) $R$-matrix, every charge order
$k\geq2$, and every homogeneous untwisted periodic chain of length $L\geq k$.
Regularity unfolds a multivariable normalized transfer matrix at any inactive edge.
Coefficients of its logarithmic derivative are then rooted connected clusters, with a
degree-$r$ coefficient supported on at most $r+1$ edges.  Subtracting the shortened
open product anchors each cluster once, so the cyclic sums of both directed endpoint
densities equal the prescribed transfer charge, including at the minimal size $L=k$.
Finally, an explicit tensor-word argument identifies the kernel of cyclic summation:
every vanishing $k$-site density is exactly a local gradient
$\mathbf 1\otimes m-m\otimes\mathbf 1$.  This upgrades equality of global charges to
the precise local equivalence in the conjecture, without a thermodynamic limit, boost
boundary terms, or removal of lower-charge admixtures.
\end{abstract}

\section{Introduction}

A regular solution of the Yang--Baxter equation supplies a commuting transfer matrix
and, by logarithmic differentiation, a hierarchy of local conserved charges.  The
hierarchy is central to quantum-integrable spin chains, while practical constructions
often use boost or ladder recurrences instead of differentiating a long monodromy
directly \cite{Tetelman1982,Links2001,TakizawaLinks2002,Bargheer2009,deLeeuw2021}.
For a general non-difference $R$-matrix, however, the choice of local density and its
relation to an algebraic representative are delicate: equal periodic sums determine
a density only modulo a local gradient, and scalar or lower-range terms cannot simply
be discarded.

Schouten and de Leeuw introduced two directed families of algebraic charge densities
in their quantum-group treatment of long-range deformations
\cite{SchoutenDeLeeuw2026}.  Their Conjecture~2.7 asserts that a difference at the two
ends of either directed block is locally equivalent to a density of the usual
logarithmic-transfer charge.  They established the first orders and the commutation of
the candidate sums with the transfer matrix, but left their identification with the
specified charge hierarchy open.  The current arXiv version at the time of writing is
v2, dated 5 June 2026; a bounded exact-title, identifier, author and conjecture-phrase
search on 14 October 2026 found no later primary proof or disproof.  This search is not
an exhaustive priority claim.

Here we prove the conjecture in its undeformed scope.  The proof has three ingredients.
First, regular permutation factors turn the two algebraic stars into the left and right
Maurer--Cartan forms of one open nearest-neighbour product.  Second, cutting a
multivariable transfer matrix at a regular factor yields an exact product over active
connected components.  Its rooted logarithmic coefficients are therefore connected
and have a sharp support bound.  Endpoint subtraction anchors each such coefficient
exactly once.  Third, a constructive tensor-word lemma shows that the kernel of cyclic
summation is precisely the space of local gradients.  Yang--Baxter integrability is
used only to equate the left and right logarithmic derivatives on the physical
homogeneous line.

The result concerns dimensionless raw derivatives at undeformed order.  It does not
extend by itself to twisted, open or inhomogeneous chains, nor to positive orders of a
long-range deformation.  In a Hermitian realization, conventional factors of $i$ and
an energy scale must still be restored model by model.

\section{Setup and statement}

Let $V$ be a finite-dimensional complex vector space and let
\begin{equation}
 R(u,v)\in\End(V\otimes V)[[u,v]]
\end{equation}
be a formal matrix power series, or the Taylor series of an analytic matrix near
$(0,0)$.  We assume the Yang--Baxter equation
\begin{equation}
 R_{12}(u_1,u_2)R_{13}(u_1,u_3)R_{23}(u_2,u_3)
 =R_{23}(u_2,u_3)R_{13}(u_1,u_3)R_{12}(u_1,u_2),
 \label{eq:ybe}
\end{equation}
regularity $R(u,u)=P$, and the fixed braiding normalization
\begin{equation}
 R_{12}(u,v)R_{21}(v,u)=\mathbf 1.
 \label{eq:braidunitarity}
\end{equation}
Equation~\eqref{eq:braidunitarity} is not Hilbert-space unitarity.  No difference-form,
Hermiticity or diagonalizability hypothesis will be used.

On a homogeneous periodic chain of length $L$, define
\begin{equation}
 t_L(u)=\tr_a\!\left[R_{aL}(u,0)\cdots R_{a1}(u,0)\right],
 \qquad U_L=t_L(0),
 \label{eq:transfer}
\end{equation}
and the raw logarithmic-transfer charges
\begin{equation}
 Q_k^{(L)}=
 \left.\partial_u^{k-2}\bigl(t_L(u)^{-1}t_L'(u)\bigr)\right|_{u=0},
 \qquad k\geq2.
 \label{eq:charges}
\end{equation}
Regularity makes $U_L$ the invertible one-site translation.  There is no factorial
division in \eqref{eq:charges}.

For ordered disjoint blocks $A=(1,\ldots,N)$ and $B=(N+1,\ldots,M)$, use the directed
product convention of Ref.~\cite{SchoutenDeLeeuw2026},
\begin{equation}
 R_{AB}=\prod_{i=1}^{N}\ \prod_{j=M}^{N+1}R_{ij}(u_i,u_j),
 \qquad D_A=\sum_{i\in A}\partial_{u_i},\quad
 D_B=\sum_{j\in B}\partial_{u_j},
\end{equation}
where the displayed order is the operator order.  The algebraic densities are
\begin{align}
 \cQ_{AB}^{(k)}&=D_A^{k-2}\left(R_{AB}^{-1}D_A R_{AB}\right),\\
 \tcQ_{AB}^{(k)}&=-D_B^{k-2}\left(R_{AB}^{-1}D_B R_{AB}\right).
 \label{eq:algebraicdensities}
\end{align}
All spectral parameters are set to zero after differentiation.  Empty-block terms are
zero.  Two $k$-site densities are locally equivalent, written $x\sim y$, when
\begin{equation}
 y-x=\mathbf 1\otimes m-m\otimes\mathbf 1
 \label{eq:gradient}
\end{equation}
for some $m\in\End(V^{\otimes(k-1)})$.

\begin{theorem}[Endpoint-density identity]
\label{thm:main}
Under \eqref{eq:ybe}--\eqref{eq:braidunitarity}, for every $k\geq2$ the two endpoint
differences
\begin{align}
 a_k={}&\left.
 \Bigl(\cQ_{1(2,\ldots,k)}^{(k)}-
       \cQ_{2(3,\ldots,k)}^{(k)}\Bigr)\right|_{\bm u=0},
 \label{eq:aksource}\\
 b_k={}&\left.
 \Bigl(\tcQ_{(1,\ldots,k-1)k}^{(k)}-
       \tcQ_{(1,\ldots,k-2)k-1}^{(k)}\Bigr)\right|_{\bm u=0}
 \label{eq:bksource}
\end{align}
are locally equivalent.  On every homogeneous untwisted periodic chain $L\geq k$,
each cyclic sum is exactly $Q_k^{(L)}$.  Consequently, both are equivalent in the
sense of \eqref{eq:gradient} to any $k$-site density of the charge
\eqref{eq:charges}.
\end{theorem}

\section{The directed blocks as open logarithmic derivatives}

Put
\begin{equation}
 C(u)=P R(u,0),\qquad C(0)=\mathbf 1,
\end{equation}
and, on sites $1,\ldots,n$, define
\begin{equation}
 G_n(u)=C_{n-1,n}(u)\cdots C_{12}(u),\qquad G_1=\mathbf 1,
 \label{eq:Gdef}
\end{equation}
together with the left and right logarithmic derivatives
\begin{equation}
 F_n=G_n^{-1}G_n',\qquad B_n=G_n'G_n^{-1}.
 \label{eq:FBdef}
\end{equation}
Superscripts on $F_n$ or $B_n$ will indicate a shifted interval of sites.

\begin{proposition}[Regular-factor reduction]
\label{prop:starreduction}
For $r=k-2$, the source expressions \eqref{eq:aksource}--\eqref{eq:bksource} are
\begin{align}
 a_k&=\left.\partial_u^r
 \left(F_k-F_{k-1}^{[2,k]}\right)\right|_{u=0},
 \label{eq:akopen}\\
 b_k&=\left.\partial_u^r
 \left(B_k-B_{k-1}^{[1,k-1]}\right)\right|_{u=0}.
 \label{eq:bkopen}
\end{align}
\end{proposition}

\begin{proof}
Let
\begin{equation}
 S_n(u)=R_{1n}(u,0)\cdots R_{12}(u,0),\qquad
 W_n=P_{1n}\cdots P_{12}.
\end{equation}
Moving each regular flip to the left gives
\begin{equation}
 W_n^{-1}S_n(u)=G_n(u).
 \label{eq:leftstar}
\end{equation}
This follows inductively: after extracting $P_{1n}$, conjugation by
$W_{n-1}$ carries $C_{1n}$ to $C_{n-1,n}$ without reversing its tensor slots.
Since $W_n$ is independent of $u$, $S_n^{-1}S_n'=F_n$.  Taking the full and
shortened stars yields \eqref{eq:akopen}.

For the opposite orientation, \eqref{eq:braidunitarity} gives, in fixed tensor slots,
\begin{equation}
 R_{ij}(0,u)=R_{ji}(u,0)^{-1}=P_{ij}C_{ij}(u)^{-1}.
 \label{eq:oppositeR}
\end{equation}
If $T_n(u)=R_{1n}(0,u)\cdots R_{n-1,n}(0,u)$, moving the flips through the
factors to their right gives
\begin{equation}
 T_n(0)^{-1}T_n(u)=
 C_{12}(u)^{-1}\cdots C_{n-1,n}(u)^{-1}=G_n(u)^{-1}.
\end{equation}
Therefore
\begin{equation}
 -T_n^{-1}T_n'=-G_n(G_n^{-1})'=G_n'G_n^{-1}=B_n,
\end{equation}
which retains the single minus sign in \eqref{eq:algebraicdensities}.  Full minus
shortened blocks now give \eqref{eq:bkopen}.  No difference property was used.
\end{proof}

\section{Cutting the multivariable transfer}

Introduce independent commuting bookkeeping variables $z_1,\ldots,z_L$ and the
normalized multivariable transfer
\begin{equation}
 X_L(\bm z)=U_L^{-1}\tr_a\!\left[
 R_{aL}(z_L,0)\cdots R_{a1}(z_1,0)\right].
 \label{eq:multix}
\end{equation}
This is only a generating series.  We do not assert that it is a new commuting
inhomogeneous family.  On the diagonal $z_1=\cdots=z_L=u$,
$X_L=U_L^{-1}t_L(u)$.  Associate $z_j$ with the directed edge $(j-1,j)$, with
indices understood cyclically.

\begin{lemma}[Exact cut and component factorization]
\label{lem:cut}
If $S$ is a proper subset of the $L$ directed edges and all variables outside $S$
are set to zero, then
\begin{equation}
 X_L\big|_S=\prod_{I\in\pi_0(S)}G_I,
 \label{eq:components}
\end{equation}
where $I$ runs over the connected interval components of $S$ and $G_I$ is the
corresponding open product \eqref{eq:Gdef}.  Products belonging to different
components commute.
\end{lemma}

\begin{proof}
Choose an inactive edge, say $z_1=0$.  Regularity inserts $R_{a1}(0,0)=P_{a1}$.
The elementary partial-trace identity
\begin{equation}
 \tr_a(A_aP_{a1})=A_1
\end{equation}
replaces the auxiliary slot by site $1$.  Equation~\eqref{eq:leftstar}, with an
independent variable in each factor, then gives
\begin{equation}
 X_L(0,z_2,\ldots,z_L)
 =C_{L-1,L}(z_L)\cdots C_{12}(z_2).
 \label{eq:onecut}
\end{equation}
A cut at any other inactive edge follows by cyclically rotating the monodromy before
the partial trace.  This use of partial-trace cyclicity is valid because physical
matrix entries of the rotated factor commute with those on all other sites.  Additional
inactive edges split \eqref{eq:onecut} into open products on disjoint site intervals.
Those factors have disjoint support and hence commute, proving
\eqref{eq:components} and independence of the chosen cut.
\end{proof}

Define the rooted logarithmic derivatives
\begin{equation}
 Y_j=X_L^{-1}\partial_{z_j}X_L,
 \qquad Z_j=(\partial_{z_j}X_L)X_L^{-1}.
 \label{eq:YZ}
\end{equation}
For a coefficient of the monomial $\bm z^{\bm\alpha}$ in either series, its rooted
edge support is
\begin{equation}
 S(\bm\alpha,j)=\{i:\alpha_i>0\}\cup\{j\}.
 \label{eq:rootedsupport}
\end{equation}
The root $j$ is included even if $\alpha_j=0$, since differentiation may have removed
its only power.

\begin{proposition}[Rooted connectedness and range]
\label{prop:rooted}
A degree-$r$ coefficient in $Y_j$ or $Z_j$ has rooted support on at most $r+1$
edges.  If $r\leq L-2$, every nonzero coefficient has connected rooted support.
\end{proposition}

\begin{proof}
The cardinality bound follows immediately from \eqref{eq:rootedsupport}.  It makes
the support proper when $r\leq L-2$, so Lemma~\ref{lem:cut} computes the coefficient
exactly after all other variables are set to zero.  If the support is disconnected,
let $I$ be the component containing the root.  Differentiating
\eqref{eq:components} and multiplying by the product inverse gives
\begin{equation}
 (X_L^{-1}\partial_{z_j}X_L)|_S=G_I^{-1}\partial_{z_j}G_I,
 \qquad
 ((\partial_{z_j}X_L)X_L^{-1})|_S=(\partial_{z_j}G_I)G_I^{-1}.
\end{equation}
Neither expression contains a variable belonging to another component.  Since every
non-root edge in $S$ has a positive exponent, the proposed disconnected coefficient
must vanish.
\end{proof}

\section{Endpoint anchoring and the global charge}

Let $\tau$ translate a local operator by one site and write the cyclic sum of a
$k$-site density $d$ on length $L$ as
\begin{equation}
 \cyc_L[d]=\sum_{n=0}^{L-1}\tau^n(d).
\end{equation}

\begin{proposition}[Each connected coefficient is counted once]
\label{prop:anchoring}
For $r=k-2$ and $L\geq k$,
\begin{align}
 \cyc_L[a_k]&=
 \left.\partial_u^r\bigl(X_L(u)^{-1}X_L'(u)\bigr)\right|_{u=0},
 \label{eq:leftsum}\\
 \cyc_L[b_k]&=
 \left.\partial_u^r\bigl(X_L'(u)X_L(u)^{-1}\bigr)\right|_{u=0}.
 \label{eq:rightsum}
\end{align}
\end{proposition}

\begin{proof}
By Proposition~\ref{prop:rooted}, a surviving degree-$r$ rooted coefficient occupies
a connected interval of at most $r+1=k-1$ edges.  It is therefore nonwrapping even at
the minimal length $L=k$.  In the left form, subtracting
$F_{k-1}^{[2,k]}$ from $F_k$ cancels exactly those rooted coefficients whose support
omits the first edge.  This remains true when the removed edge is the root and its
exponent after differentiation is zero: such a term has no counterpart in the shortened
product.  Thus every remaining interval is anchored at its unique leftmost edge.
Cyclic translation counts every connected rooted coefficient exactly once.

For the right form, subtracting $B_{k-1}^{[1,k-1]}$ instead anchors the unique
rightmost edge, and the same argument applies.  Raw differentiation multiplies every
degree-$r$ coefficient by the same $r!$ on both sides.  The case $k=2$ is included:
the rooted support consists only of its root and $G_1'=0$.
\end{proof}

The Yang--Baxter equation now enters.  Its standard train argument gives
$[t_L(u),t_L(v)]=0$.  Hence $t_L$, $t_L'$ and $U_L=t_L(0)$ commute, and on the
homogeneous line
\begin{equation}
 X_L^{-1}X_L'=t_L^{-1}t_L'
 =X_L'X_L^{-1}.
 \label{eq:formscoincide}
\end{equation}
Combining Proposition~\ref{prop:anchoring} with \eqref{eq:formscoincide} proves
\begin{equation}
 \cyc_L[a_k]=Q_k^{(L)}=\cyc_L[b_k]
 \qquad (L\geq k).
 \label{eq:globalsums}
\end{equation}
No infinite-chain limit or finite periodic boost is involved, and scalar and lower-range
coefficients remain in the calculation.

\section{The exact kernel of cyclic summation}

Global equality is not yet the local equivalence \eqref{eq:gradient}.  The required
upgrade follows from a finite-dimensional algebraic lemma.

\begin{lemma}[Cyclic kernel]
\label{lem:kernel}
Let $D\in\End(V^{\otimes k})$.  If $\cyc_L[D]=0$ for one $L>2k$, then there exists
$M\in\End(V^{\otimes(k-1)})$ such that
\begin{equation}
 D=\mathbf 1\otimes M-M\otimes\mathbf 1.
 \label{eq:kernelgradient}
\end{equation}
Conversely, every expression of the form \eqref{eq:kernelgradient} has zero cyclic sum
for every $L\geq k$.
\end{lemma}

\begin{proof}
Choose a basis of $\End(V)$ consisting of $\mathbf 1$ and trace-zero matrices
$E_\alpha$, and expand $D$ in the induced tensor-word basis.  For any word other than
the all-identity word, delete its leading and trailing identities.  The remaining
\emph{core} $w$ has length $m\leq k$ and nonidentity first and last letters.  Let
$c_a$ be the coefficient of the embedding
\begin{equation}
 \mathbf 1^{\otimes a}\otimes w\otimes
 \mathbf 1^{\otimes(k-m-a)},\qquad 0\leq a\leq k-m.
\end{equation}
On a ring with $L>2k$, the exterior identity run uniquely identifies the core and its
cyclic placement, so distinct cores cannot cancel one another.  Vanishing of the cyclic
sum is therefore equivalent to a zero scalar coefficient and
\begin{equation}
 \sum_{a=0}^{k-m}c_a=0
 \label{eq:coresum}
\end{equation}
for every core.

For $m<k$, put the same core into a $(k-1)$-site word and assign coefficients
\begin{equation}
 d_a=-\sum_{b=0}^{a}c_b,\qquad 0\leq a\leq k-m-1.
 \label{eq:partialsums}
\end{equation}
Then $\mathbf 1\otimes M-M\otimes\mathbf 1$ has coefficients
$-d_0=c_0$, $d_{a-1}-d_a=c_a$ in the interior, and
$d_{k-m-1}=c_{k-m}$ at the final embedding, where the last equality is
\eqref{eq:coresum}.  Summing these contributions over all cores constructs $M$.
For $m=k$, \eqref{eq:coresum} forces the sole coefficient to vanish.  If
$\dim V=1$, only the scalar word exists and is likewise forced to zero.  The converse
is immediate by telescoping around the ring.
\end{proof}

Apply Lemma~\ref{lem:kernel} to $D=a_k-b_k$, using \eqref{eq:globalsums} at any
$L>2k$.  This proves $a_k\sim b_k$.  If $q_k$ is any $k$-site representative whose
cyclic sum is $Q_k^{(L)}$ for all $L\geq k$, the same lemma applied to $a_k-q_k$ and
$b_k-q_k$ proves both equivalences claimed in Theorem~\ref{thm:main}.  Alternatively,
$a_k$ itself supplies a canonical representative, so no independently normalized
density need be assumed.  This completes the proof of Theorem~\ref{thm:main}.

\section{Exact finite checks and convention control}

The theorem is analytic; the following computations test conventions rather than supply
its universal quantifiers.  Two independent implementations transcribed the transfer
matrix and both directed block products using exact integer or rational arithmetic.
Both implementations were developed with OpenAI Codex assistance; their independence
concerns the input reconstruction and matrix representation, not independent human
authorship.  Generated code was checked by complete exact residuals and a separate
fresh AI reconstruction.  The
published spin-$\tfrac12$ XXX input
\begin{equation}
 R_{12}(u,v)=\frac{(u-v)\mathbf 1-P_{12}}{u-v-1}
\end{equation}
reproduces the known $Q_2,Q_3,Q_4$ densities after the explicit normalization map to
Ref.~\cite{Bargheer2009}.  A nonlinear reparametrization $f(u)=u+u^2$ retains the
chain-rule admixtures $Q_4[f]=Q_4+6Q_3$ and
$Q_5[f]=Q_5+12Q_4+12Q_3$; both endpoint formulas solve the complete local-gradient
system at $k=4,5$.  An admissible spectral-gauge transform provides a separate
non-difference control through $k=6$, including $L=k=6$.

The recorded exact checks are as follows; ``zero'' means entrywise exact zero, with no
floating-point tolerance.
\begin{itemize}
\item For XXX at $k\leq4$ and $L=4,5,6$, independent transfer and block
constructions have zero residual.
\item For $f(u)=u+u^2$ at $k=4$ with $L=4,5,6$ and at $k=5$ with $L=5,6$,
the independent transfers agree exactly; both complete local-gradient systems
also have zero residual, and hence the endpoint sums equal the same charge.
\item For spectral-gauge XXX at $k=4,5,6$ and $L=6,7$, the sums, gradients and
independent transfer comparison have zero residual.
\item For generic non-Yang--Baxter controls through $k=5$, the cut and
connected-support identities have zero residual.
\end{itemize}

The generic controls in the last row are deliberately outside the physical target and
are not counterexamples or confirmations of the Yang--Baxter theorem.  They isolate the
purely algebraic cut and support steps.  Reproduction commands, scripts, exact outputs
and the Python version are specified in the accompanying reproducibility record.

\section{Physical meaning and limitations}

The theorem identifies the algebraic densities of Ref.~\cite{SchoutenDeLeeuw2026}
with the particular hierarchy selected by the logarithmic derivative
\eqref{eq:charges}; they are not merely another mutually commuting family.  In
particular, no untracked scalar or lower-charge combination remains after passage to the
exact local-gradient quotient.  This supplies the missing undeformed identification
used conditionally in that work and gives finite-support representatives at every order.

For a Hermitian XXX realization one may write $H=JQ_2$ with real $J$ carrying energy
units.  Odd raw derivatives can require the customary factor of $i$ to become Hermitian.
The general algebraic $R$ in Theorem~\ref{thm:main} need not be a unitary quantum gate,
and the theorem does not predict a new excitation, phase or experimental signal.  Its
scope is the homogeneous, untwisted, periodic, undeformed chain.  Open boundaries,
twists, inhomogeneous evaluation parameters and positive long-range deformation orders
can introduce additional endpoint or wrapping terms and are not covered.

\section{Conclusion}

Regularity turns an inactive transfer factor into an exact cut.  This makes the
logarithmic-transfer coefficients rooted connected clusters and supplies a support bound
strong enough to retain the minimal chain length $L=k$.  The two directed endpoint
subtractions anchor those clusters from opposite ends, while Yang--Baxter commutativity
identifies the resulting left and right logarithmic derivatives.  A constructive cyclic
kernel lemma then yields the precise local gradient.  Together these steps prove both
orientations of Conjecture~2.7 of Ref.~\cite{SchoutenDeLeeuw2026} for general regular
braiding-normalized $R$-matrices and all charge orders.

\section*{Data and code availability}

The exact-arithmetic scripts, machine-readable outputs, primary-source locators and
regeneration commands accompanying this local draft are described in
\texttt{REPRODUCIBILITY.md}.  They test finite cases and conventions; the all-order
claim rests on the proof above.  No third-party article PDF is part of the proposed
submission bundle.

\section*{Acknowledgements and AI assistance disclosure}

OpenAI Codex was used to support literature retrieval and organization, proof
development, exact-arithmetic code, an independent AI audit, and manuscript drafting.
AI-assisted outputs were checked by exact replays and by a separate fresh AI scientific
review, but those checks are not external peer review.  This is a local draft: the named
author must critically review and validate all content, approve the disclosure, and take
responsibility for any version submitted to a journal.

\bibliographystyle{unsrt}
\bibliography{references}

\end{document}
